\documentclass[14pt,a4paper,oneside]{extbook}
\usepackage[utf8]{vietnam}
\usepackage{amsmath,amssymb,mathrsfs,fancyhdr,enumerate,multirow,makecell}
\usepackage{tkz-euclide,tikz-3dplot,tikz,tkz-tab}
\usetikzlibrary{arrows.meta,calc,intersections,angles,quotes,backgrounds,shapes.geometric,patterns,through,shadows,positioning}
\usepackage{yhmath}
\usepackage[top=0.7cm, bottom=0.7cm, left=1.0cm, right=1.0cm, headheight=16pt, headsep=3mm, footskip=5mm]{geometry}
\usepackage[hidelinks,unicode]{hyperref}
\usepackage{qrcode}
\usepackage[many]{tcolorbox}
\usepackage{fontawesome}

\pagestyle{fancy}
\fancyhf{}
\fancyhead[L]{\footnotesize\bfseries\color{teal!85!black}\faBook~SỔ TAY KIẾN THỨC TOÁN 4}
\fancyhead[R]{\footnotesize\bfseries\color{blue!85!black}Lớp Toán Thầy Lĩnh -- loptoanthaylinh.vn}
\fancyfoot[C]{\footnotesize\bfseries\color{gray!80!black}-- Sổ tay kiến thức: Bài 54. Phân số và phép chia số tự nhiên --}
\renewcommand{\headrulewidth}{0.4pt}

\begin{document}

% --- BANNER TIÊU ĐỀ BÀI HỌC ---
\noindent
\begin{tcolorbox}[
  enhanced,
  colback=teal!8!white,
  colframe=teal!80!black,
  arc=3mm, boxrule=1pt,
  left=3mm, right=3mm, top=2mm, bottom=2mm,
  halign=center
]
  {\Large\bfseries\color{teal!90!black}\faBookmark\ BÀI 54: PHÂN SỐ VÀ PHÉP CHIA SỐ TỰ NHIÊN}\\[1.5pt]
  {\footnotesize\itshape\color{black!75} Mối quan hệ phép chia a : b = a/b, so sánh phân số với 1 trên 1 trang A4}
\end{tcolorbox}

\vspace{-2.5mm}

\noindent
\begin{minipage}[t]{0.49\textwidth}\vspace{0pt}
\begin{tcolorbox}[
  enhanced,
  colback=blue!4!white,
  colframe=blue!75!black,
  arc=2mm, boxrule=0.8pt,
  left=2mm, right=2mm, top=1.8mm, bottom=1.8mm,
  title={\footnotesize\bfseries\color{white}\faShareAlt\ 1. Thương của phép chia số tự nhiên}
]
{\footnotesize
Thương của phép chia số tự nhiên cho số tự nhiên khác 0 luôn viết được thành \textbf{phân số}:
$$a : b = \dfrac{a}{b} \quad (b \ne 0)$$
\begin{itemize}\setlength{\itemsep}{1pt}\setlength{\leftskip}{-2.5mm}
  \item Tử số là \textbf{Số bị chia}, Mẫu số là \textbf{Số chia}.
  \item Mọi số tự nhiên đều viết được thành phân số có \textbf{mẫu số là 1}:
  $$a = \dfrac{a}{1} \quad (5 = \dfrac{5}{1},~ 100 = \dfrac{100}{1})$$
  \item Số 0 viết thành phân số: $0 = \dfrac{0}{b}$ ($b \ne 0$). Số 1 viết thành: $1 = \dfrac{a}{a}$ ($a \ne 0$).
\end{itemize}
}
\end{tcolorbox}
\end{minipage}
\hfill
\begin{minipage}[t]{0.49\textwidth}\vspace{0pt}
\begin{tcolorbox}[
  enhanced,
  colback=teal!4!white,
  colframe=teal!75!black,
  arc=2mm, boxrule=0.8pt,
  left=2mm, right=2mm, top=1.8mm, bottom=1.8mm,
  title={\footnotesize\bfseries\color{white}\faBalanceScale\ 2. So sánh phân số với 1}
]
{\footnotesize
So sánh tử số ($a$) và mẫu số ($b$):
\begin{itemize}\setlength{\itemsep}{1.5pt}\setlength{\leftskip}{-2.5mm}
  \item \textbf{Tử số bé hơn mẫu số ($a < b$):}\\
  Phân số \textbf{bé hơn 1}: $\dfrac{a}{b} < 1$ (Ví dụ: $\frac{3}{5} < 1$).
  \item \textbf{Tử số bằng mẫu số ($a = b$):}\\
  Phân số \textbf{bằng 1}: $\dfrac{a}{b} = 1$ (Ví dụ: $\frac{4}{4} = 1$).
  \item \textbf{Tử số lớn hơn mẫu số ($a > b$):}\\
  Phân số \textbf{lớn hơn 1}: $\dfrac{a}{b} > 1$ (Ví dụ: $\frac{7}{4} > 1$).
\end{itemize}
}
\end{tcolorbox}
\end{minipage}

\vspace{-2.5mm}
% --- KHỐI 3: SƠ ĐỒ TƯ DUY TIKZ ---
\noindent
\begin{tcolorbox}[
  enhanced,
  colback=orange!4!white,
  colframe=orange!80!black,
  arc=2mm, boxrule=0.8pt,
  left=2mm, right=2mm, top=1.5mm, bottom=1.5mm,
  title={\footnotesize\bfseries\color{white}\faSitemap\ SƠ ĐỒ TƯ DUY: CÁN CÂN SO SÁNH PHÂN SỐ VỚI 1}
]
\centering
\begin{tikzpicture}[
  box/.style={draw, rounded corners=2pt, align=center, font=\scriptsize\bfseries, inner sep=3pt, minimum height=0.95cm},
  arr/.style={-{Stealth[length=2.5mm]}, thick}
]
  \node[box, fill=purple!12, draw=purple!80!black, text width=3.2cm] (root) at (0, 0) {PHÂN SỐ $\dfrac{a}{b}$\\So sánh Tử \& Mẫu};

  \node[box, fill=blue!10, draw=blue!80!black, text width=3.5cm] (b1) at (5.2, 1.15) {TỬ $<$ MẪU ($a < b$)\\$\implies \mathbf{\dfrac{a}{b} < 1}$\\Ví dụ: $\frac{2}{5} < 1$};

  \node[box, fill=teal!10, draw=teal!80!black, text width=3.5cm] (b2) at (5.2, 0.0) {TỬ $=$ MẪU ($a = b$)\\$\implies \mathbf{\dfrac{a}{b} = 1}$\\Ví dụ: $\frac{6}{6} = 1$};

  \node[box, fill=orange!12, draw=orange!80!black, text width=3.5cm] (b3) at (5.2, -1.15) {TỬ $>$ MẪU ($a > b$)\\$\implies \mathbf{\dfrac{a}{b} > 1}$\\Ví dụ: $\frac{8}{3} > 1$};

  \draw[arr, blue!80!black] (root.east) -- ++(0.5, 0) |- (b1.west);
  \draw[arr, teal!80!black] (root.east) -- (b2.west);
  \draw[arr, orange!80!black] (root.east) -- ++(0.5, 0) |- (b3.west);
\end{tikzpicture}
\end{tcolorbox}

\vspace{-2.5mm}
% --- KHỐI 4: VÍ DỤ & MẸO TRÁNH BẪY ---
\noindent
\begin{minipage}[t]{0.49\textwidth}\vspace{0pt}
\begin{tcolorbox}[
  enhanced,
  colback=green!4!white,
  colframe=green!60!black,
  arc=2mm, boxrule=0.8pt,
  left=2mm, right=2mm, top=1.8mm, bottom=1.8mm,
  title={\footnotesize\bfseries\color{white}\faCheckCircle\ Ví dụ thực hành chuẩn mực}
]
{\footnotesize
\textbf{1. Viết phép chia dưới dạng phân số:}\\
$3 : 7 = \mathbf{\dfrac{3}{7}}; \quad 5 : 8 = \mathbf{\dfrac{5}{8}}; \quad 9 : 3 = \mathbf{\dfrac{9}{3}} = \mathbf{3}$.\\
\textbf{2. So sánh các phân số sau với 1:}\\
$\bullet~ \dfrac{4}{7} < 1$ (vì $4 < 7$); \quad $\bullet~ \dfrac{9}{9} = 1$ (vì tử = mẫu);\\
$\bullet~ \dfrac{11}{8} > 1$ (vì $11 > 8$).
}
\end{tcolorbox}
\end{minipage}
\hfill
\begin{minipage}[t]{0.49\textwidth}\vspace{0pt}
\begin{tcolorbox}[
  enhanced,
  colback=red!4!white,
  colframe=red!70!black,
  arc=2mm, boxrule=0.8pt,
  left=2mm, right=2mm, top=1.8mm, bottom=1.8mm,
  title={\footnotesize\bfseries\color{white}\faExclamationTriangle\ Mẹo tránh bẫy trắc nghiệm}
]
{\footnotesize
\begin{itemize}\setlength{\itemsep}{1.5pt}\setlength{\leftskip}{-2.5mm}
  \item \textbf{Bẫy viết ngược vị trí:}\\
  $3 : 5$ phải viết là $\dfrac{3}{5}$, nhiều học sinh hay viết ngược thành $\dfrac{5}{3}$ là SAI.
  \item \textbf{Bẫy số 1 dưới mẫu:}\\
  Số tự nhiên $7 = \dfrac{7}{1}$, không được viết thành $\dfrac{1}{7}$!
\end{itemize}
}
\end{tcolorbox}
\end{minipage}

\end{document}
